%HW07.tex
%
% Sixth Homework for Honors Algebra I  (415H)
% Frank Sottile
%%%%%%%%%%%%%%%%%%%%%%%%%%%%%%%%%%%%%%%%%%%%%%%%%%%%%%%%%%%%%%%%%%%%%%%
\documentclass[12pt]{article}
\usepackage{multicol,amsfonts, amssymb,  mathtools,amsmath}
\usepackage[dvipsnames]{xcolor}
\usepackage{colordvi,graphicx}
\headheight=8pt
%
\topmargin=-75pt
\textheight=720pt   \textwidth=560pt
\oddsidemargin=-60pt \evensidemargin=-60pt

\pagestyle{empty}   %  Change this for your write-up

%%%%%%%%%%%%%%%%%%%%%%%%%%%%%%%%%%%%%%%%%%%%

\newcommand{\CC}{{\mathbb C}}
\newcommand{\NN}{{\mathbb N}}
\newcommand{\QQ}{{\mathbb Q}}
\newcommand{\RR}{{\mathbb R}}
\newcommand{\ZZ}{{\mathbb Z}}

\newcommand{\Aut}{\textrm{Aut}}
\newcommand{\op}{\textrm{op}}

\newcommand{\vect}[2]{(\begin{smallmatrix}#1\\#2\end{smallmatrix})}

\def\Color#1#2{\special{color push cmyk #1}#2\special{color pop}}
%\def\Indigo#1{\Color{.42 1. 0. .49}{#1}}
\def\Indigo#1{\Color{1. .95 .05 .4}{#1}}
\def\MyViolet#1{\Color{.6 1. 0. .15}{#1}}
\def\TAMU#1{\Color{.15 1. .39 .49}{#1}}

\newcommand{\barsl}{\noindent\begin{minipage}[t]{590pt}
\Indigo{\rule{590pt}{1.2pt}}\vspace{-5.7mm}\\
\MyViolet{\rule{590pt}{1.2pt}}\vspace{-5.7mm}\\
\Blue{\rule{590pt}{1.2pt}}\vspace{-5.7mm}\\
\Green{\rule{590pt}{1.2pt}}\vspace{-5.7mm}\\
\Yellow{\rule{590pt}{1.2pt}}\vspace{-5.7mm}\\
\Orange{\rule{590pt}{1.2pt}}\vspace{-5.7mm}\\
\Red{\rule{590pt}{1.2pt}}\bigskip
\end{minipage}}


\newcommand{\barsn}{\noindent\begin{minipage}[t]{560pt}
\Indigo{\rule{560pt}{1.1pt}}\vspace{-4.5mm}\\
\MyViolet{\rule{560pt}{1.1pt}}\vspace{-4.5mm}\\
\Blue{\rule{560pt}{1.1pt}}\vspace{-4.5mm}\\
\Green{\rule{560pt}{1.1pt}}\vspace{-4.5mm}\\
\Yellow{\rule{560pt}{1.1pt}}\vspace{-4.5mm}\\
\Orange{\rule{560pt}{1.1pt}}\vspace{-4.5mm}\\
\Red{\rule{560pt}{1.1pt}}\bigskip
\end{minipage}}

\def\demph#1{\TAMU{{\sl #1}}}
\def\defcolor#1{\TAMU{#1}}

\begin{document}
\LARGE 
\noindent
Algebra 415H \ \ Autumn 2026\vspace{2pt}\\
{\color{magenta}Your Name}\vspace{2pt}\\
\Large 8 October 2026 \hfill
\sf
 Seventh Homework\makebox[40pt][l]{\ }
\normalsize\vspace{10pt}

\noindent
Write your answers neatly, in complete sentences.  
I highly recommend recopying your work before handing it in.
Correct and crisp proofs are greatly appreciated; oftentimes your work can be shortened and made clearer.

\noindent
I also recommend talking with other students, and I am able to give hints to public posts in Piazza.

\noindent
\barsn

\noindent\TAMU{{\sf Hand in at the start of class, Thursday 15 October:}} 


\begin{enumerate}
\setcounter{enumi}{45}

%%%%%%%%%%%%%%%%%%%%%%%%%%%%%%%%%%%%%%%%%%%%%%%%%%%%%%%%%%%%%%%%%%%%%%%%%%%%%%%%%%%%%%%%%%%%%%%%%%%%
\item
  Suppose that $G$ is a finite group of order $n$.
  Prove that every element $a\in G$ satisfies $a^n=e$.
  
%%%%%%%%%%%%%%%%%%%%%%%%%%%%%%%%%%%%%%%%%%%%%%%%%%%%%%%%%%%%%%%%%%%%%%%%%%%%%%%%%%%%%%%%%%%%%%%%%%%%

%%%%%%%%%%%%%%%%%%%%%%%%%%%%%%%%%%%%%%%%%%%%%%%%%%%%%%%%%%%%%%%%%%%%%%%%%%%%%%%%%%%%%%%%%%%%%%%%%%%%
\item
  Let $H$ and $K$ be subgroups of a group $G$.
  Define the relation \defcolor{$\sim$} on $G$ by \defcolor{$a\sim b$} if and only if there are $h\in H$ and $k\in K$ such that
  $a=hbk$.
  \begin{enumerate}
    \item
      Prove that $\sim$ is an equivalence relation.
    \item
      Describe the equivalence class $[a]$ in a succinct form.
    \item
      Supppose that $G=S_3$ and $H=K=\{e, (1,2)\}$.
      Find all $\sim$ equivalence classes.
  \end{enumerate}
%%%%%%%%%%%%%%%%%%%%%%%%%%%%%%%%%%%%%%%%%%%%%%%%%%%%%%%%%%%%%%%%%%%%%%%%%%%%%%%%%%%%%%%%%%%%%%%%%%%%

%%%%%%%%%%%%%%%%%%%%%%%%%%%%%%%%%%%%%%%%%%%%%%%%%%%%%%%%%%%%%%%%%%%%%%%%%%%%%%%%%%%%%%%%%%%%%%%%%%%%
\item
  Let $G$ be a group.  Earlier we defined the \demph{center} of $G$ to be the subgroup
  \[
     \defcolor{C(G)}\ \vcentcolon=\ \{ g\in G \mid gh=hg \mbox{ for all }h\in G\}\,.
  \]
  \begin{enumerate}
    \item
      What is the center of the group $D_4=\langle \sigma,\rho\rangle$, where $\sigma$ is a reflection and $\rho$ is the
      rotation through an angle of $\pi/2$?

    \item
      Prove that the center of a group is a normal subgroup.
      
    \item
      Identify $D_4/C(D_4)$.
      
    \item
      Prove that if the quotient group $G/C(G)$ is cyclic, then $G$ is abelian.
  \end{enumerate}  
%%%%%%%%%%%%%%%%%%%%%%%%%%%%%%%%%%%%%%%%%%%%%%%%%%%%%%%%%%%%%%%%%%%%%%%%%%%%%%%%%%%%%%%%%%%%%%%%%%%%

%%%%%%%%%%%%%%%%%%%%%%%%%%%%%%%%%%%%%%%%%%%%%%%%%%%%%%%%%%%%%%%%%%%%%%%%%%%%%%%%%%%%%%%%%%%%%%%%%%%%
\item
  Let $H$ be a subgroup of a group $G$ and $a,b\in G$.
  For each of the following, prove or give a counter example.

  {\quad}\quad  \makebox[3in][l]{(a) If $aH=bH$, then $Ha=Hb$.}
  \ 
   (b) If $a^{-1}H=b^{-1}H$, then $Ha=Hb$.

%%%%%%%%%%%%%%%%%%%%%%%%%%%%%%%%%%%%%%%%%%%%%%%%%%%%%%%%%%%%%%%%%%%%%%%%%%%%%%%%%%%%%%%%%%%%%%%%%%%%

%%%%%%%%%%%%%%%%%%%%%%%%%%%%%%%%%%%%%%%%%%%%%%%%%%%%%%%%%%%%%%%%%%%%%%%%%%%%%%%%%%%%%%%%%%%%%%%%%%%%
\item
  Show that a finite group cannot be written as a union of two of its subgroups.

  What if we replace ``two'' by ``three''?   (Prove or give a counter example.)

%%%%%%%%%%%%%%%%%%%%%%%%%%%%%%%%%%%%%%%%%%%%%%%%%%%%%%%%%%%%%%%%%%%%%%%%%%%%%%%%%%%%%%%%%%%%%%%%%%%%



%%%%%%%%%%%%%%%%%%%%%%%%%%%%%%%%%%%%%%%%%%%%%%%%%%%%%%%%%%%%%%%%%%%%%%%%%%%%%%%%%%%%%%%%%%%%%%%%%%%%
\item
  Let $G,H$ be groups.
  Define the function $\varphi\colon G\times H \to G$ given by $(g,h)\mapsto g$.
  Show that this is a homomorphism.
  Determine its kernel.
  
%%%%%%%%%%%%%%%%%%%%%%%%%%%%%%%%%%%%%%%%%%%%%%%%%%%%%%%%%%%%%%%%%%%%%%%%%%%%%%%%%%%%%%%%%%%%%%%%%%%%

%%%%%%%%%%%%%%%%%%%%%%%%%%%%%%%%%%%%%%%%%%%%%%%%%%%%%%%%%%%%%%%%%%%%%%%%%%%%%%%%%%%%%%%%%%%%%%%%%%%%
\item
  Two elements $a$ abnd $b$ of a group $G$ are \demph{conjugate} if there is an element $g\in G$ such that $a=gbg^{-1}$
  (this is an equivalence relation whose equivalence classes are called conjugacy classes).
  Determine all the conjugacy classes in $S_3$ and in $S_4$.
  There is a hard way to figure this out and an easy way.

  If you understand the easy way, or think about your answer, you should be able to conjecture (and prove) what are the
  conjugacy classes in \demph{any} given symmetric group $S_n$, which I ask for you to do.

%%%%%%%%%%%%%%%%%%%%%%%%%%%%%%%%%%%%%%%%%%%%%%%%%%%%%%%%%%%%%%%%%%%%%%%%%%%%%%%%%%%%%%%%%%%%%%%%%%%%


%%%%%%%%%%%%%%%%%%%%%%%%%%%%%%%%%%%%%%%%%%%%%%%%%%%%%%%%%%%%%%%%%%%%%%%%%%%%%%%%%%%%%%%%%%%%%%%%%%%%
\item
  Let $H,K$ be subgroups of a group $G$.
  We say that \demph{$H$ is conjugate to $K$} if there exists $g\in G$ such that $H = gKg^{-1}$.
  Show that conjugacy is an equivalence relation on the set of subgroups of $G$.

%%%%%%%%%%%%%%%%%%%%%%%%%%%%%%%%%%%%%%%%%%%%%%%%%%%%%%%%%%%%%%%%%%%%%%%%%%%%%%%%%%%%%%%%%%%%%%%%%%%%


%%%%%%%%%%%%%%%%%%%%%%%%%%%%%%%%%%%%%%%%%%%%%%%%%%%%%%%%%%%%%%%%%%%%%%%%%%%%%%%%%%%%%%%%%%%%%%%%%%%%
\item
  Let $G$ be a group.
  Let \defcolor{$\Aut(G)$} be the set of automorphisms of $G$.
  \begin{enumerate}
    \item
       Show that $\Aut(G)$ is a group under function composition.

     \item
       Show that the set of inner automorphisms forms a subgroup of  $\Aut(G)$, and that it is normal.
  \end{enumerate}  
%%%%%%%%%%%%%%%%%%%%%%%%%%%%%%%%%%%%%%%%%%%%%%%%%%%%%%%%%%%%%%%%%%%%%%%%%%%%%%%%%%%%%%%%%%%%%%%%%%%%




\end{enumerate}
%%%%%%%%%%%%%%%%%%%%%%%%%%%%%%%%%%%%%%%%%%%%%%%%%%%%%%%%%%%%%%%%%%%%%%%%%%%%%%%%%%%%%%%%%%%%%%%%%%%%

\end{document}


%%%%%%%%%%%%%%%%%%%%%%%%%%%%%%%%%%%%%%%%%%%%%%%%%%%%%%%%%%%%%%%%%%%%%%%%%%%%%%%%%%%%%%%%%%%%%%%%%%%%
\item

%%%%%%%%%%%%%%%%%%%%%%%%%%%%%%%%%%%%%%%%%%%%%%%%%%%%%%%%%%%%%%%%%%%%%%%%%%%%%%%%%%%%%%%%%%%%%%%%%%%%
